\textsc{Minimum Vertex Cover} (MVC) es NP-difícil, y bajo la conjetura de juegos únicos no admite aproximaciones mejores que $2$. Presentamos su marco teórico e histórico y estudiamos tres algoritmos de paradigmas distintos: el voraz de grado máximo de Weigt y Hartmann, la programación dinámica sobre descomposiciones en árbol y el conjunto independiente maximal aleatorio de Veldt, que es una $2$-aproximación en esperanza en tiempo lineal. Para cada uno explicamos su enfoque, damos su pseudocódigo, demostramos su complejidad y describimos en qué tipos de grafos funciona. De la comparación de los tres derivamos \hyvc, un algoritmo híbrido que reduce el grafo con reglas exactas, lo divide en componentes, resuelve por programación dinámica las componentes angostas y ramifica sobre el resto con cotas superiores de los dos heurísticos y la cota LP de Nemhauser y Trotter. \hyvc{} siempre devuelve una cubierta válida, certifica optimalidad cuando termina y corre en $O^*(1.4656^n)$. Implementamos todo en C++ y lo evaluamos en 435 instancias públicas o reproducibles: las 200 del PACE~2019, las 200 del PACE~2017 (con ancho de árbol conocido) y 35 sintéticas. En 10\,s, \hyvc{} resuelve 253 de ellas, frente a 158 de la programación dinámica y 83 del divide y vencerás exponencial. En el PACE~2019 queda muy por detrás del ganador, \textsc{WeGotYouCovered} (30 frente a 155 instancias), y mostramos que la diferencia se debe a reglas de reducción y cotas inferiores que ninguno de los tres algoritmos aporta.
